% Terminal apparatus of the concise study: the scope note and the References
% for the sources this companion actually uses. The revision timestamp and the
% rights colophon follow from the shared generation record.

\section{Appendix: Scope and Qualifications}

{\footnotesize
\textbf{Edition, formulary and occurrence.} This concise study treats the Mass
\latin{Dominica decima octava post Pentecosten}, II classis, of the
\work{Missale Romanum}, editio typica, Typis Polyglottis Vaticanis 1962,
pp.~410--411, nos.~1669--1678, in the universal 1962 calendar and no particular
calendar. Its occurrence on 27~September 2026 rests on the Missal's own calendar
and general rubrics (RG 111~b, 127~b; RGMR 427, 434~b, 494~b), and agrees with
the dated 2026 Ordos of the Fraternity of St Peter and of the Institute of
Christ the King in France, cited for the date only; the particular solemnity
each prints for its own institute is no part of the universal formulary.

\textbf{Relation to the expansive study.} Every claim here is drawn from the
canonical study of this same formulary, \textit{The Eighteenth Sunday after
Pentecost}, which prints the appointed Latin with its public-domain English,
treats each element in its own setting, and develops the three readings at
length, each with its own four senses; nothing evidence-dependent is added here.
The four rows on the first page orient a reader and do not replace those
readings' own senses. The map and the dossier reprint the study's, with the
dates, relation labels, alternatives and unresolved states exactly as the
generated chronology record gives them. The joins between elements are the
editor's; each Father and later writer is cited only for what he said of his
own passage, and the liturgical commentators for the elements of this Mass
they expound.

\textbf{Text, English and rights.} Every Latin form of the ten elements was read
for the study on the page images of the CMAA facsimile of the typical edition
and compared with the Pustet \work{Missale Romanum} of Ratisbon, 1862, in which
each element's wording already stands; the collation is in
\texttt{propers/verified.md}. The Latin source of the Introit's and the
Offertory's distinctive wording, and the psalter behind the Communion's
\latin{in aula sancta}, were not identified. The scriptural English is the
Douay--Rheims in Challoner's revision at the canonical verses, so that it does
not render the Missal's Latin where a chant adapts its source; the English of
the Collect is the anonymous translation printed by Cummiskey in Philadelphia
in 1861. Neither is an approved liturgical translation of the 1962 Missal, and
English renderings of patristic Latin not drawn from a named translation are
glosses of the Latin beside them. The Latin of the formulary rests on its
public-domain antecedent, the Pustet Missal (17 U.S.C. 103(b)); the
Douay--Rheims, the 1861 English, the Fathers in Migne and Zingerle, and the
nineteenth- and early twentieth-century English of Chrysostom, Augustine,
Bellarmine, Schuster and the continuation of \work{The Liturgical Year} are in
the public domain; the Corpus Thomisticum text of Aquinas on First Corinthians
is quoted briefly from a modern presentation; the \work{New American Bible}
introduction and the Ordos are cited, not reproduced.

\textbf{Reception and its limits.} The witnesses named here are those the
study's sweep read at the loci below. Chrysostom on Matthew and Augustine on the
psalms and the harmony of the Gospels are quoted from the Nicene and
Post-Nicene English, with Augustine's Latin read in Migne at the loci quoted;
Chrysostom on First Corinthians and Hebrews from New Advent's delivery of the
same translation. The Greek of Chrysostom on Ps~121, of Theodoret and of the
expositions under St Athanasius's name was read in Migne's columns. Bellarmine
was read only in O'Sullivan's abridged English of 1866, in a web transcription;
Cornelius a Lapide in optical text layers of the Antwerp printings, and he is
summarised rather than quoted; St Ambrose and St Bede on Luke in transcriptions
of Migne, not on page images; St Hilary on Matthew in the Maurist text, not
its modern critical edition. No patristic commentary on Ecclus 36:18 or on the
Offertory's cited verses of Exodus was located. The continuation of \work{The
Liturgical Year} is Dom Lucien Fromage's, by the library's standing registry,
and is never cited as Guéranger's. The sacramentary witnesses are Wilson's
editions, the chant witnesses gregorien.info's report of Hesbert's
\work{Sextuplex}, which was not itself opened, and the lectionaries Morin's; no
manuscript and no critical edition of
the sacramentaries was consulted, and none of these witnesses shows when this
Gospel first stood beside this Offertory.

\textbf{Chronology and records.} The Date cells on the second page come from
the Scripture chronology corpus through the record generated for this
formulary. That record gives the three psalms only the modern critical boundary
for the Psalter, from the introduction to the Psalms in the \work{New American
Bible, Revised Edition} (summarized, not quoted), and holds no traditional date
for them. For the Gospel and the Epistle it holds the traditional dates the
\work{Catholic Encyclopedia} reports; for the Communion's psalm a superscription
setting and a prophetic referent; for the Offertory a relative date within the
narrative; and for the Introit a separate answer at each of its two loci, so
that the element has no single date. The Gospel's narrated event is undated.
The search, its loci, the disagreements and the negative results are in
\texttt{research/scope.md}, the reasoning behind the three readings in
\texttt{research/interpretations.md}, and this companion's production record in
\texttt{research/production-review.md}.
\par}

\section*{References}
\runninghead{References}

{\scriptsize
\subsubsection*{Liturgical books and Bibles}
\begin{itemize}\setlength{\itemsep}{0.15em}\setlength{\parskip}{0pt}
\item \work{Missale Romanum ex decreto Sacrosancti Concilii Tridentini restitutum, Summorum Pontificum cura recognitum}, editio typica (Typis Polyglottis Vaticanis, 1962), CMAA facsimile: \work{Rubricae generales} 111, 127; \work{Rubricae generales Missalis} 427, 434, 494; pp.~410--411 (nos.~1669--1678).
\item \work{Missale Romanum} (Ratisbon: Pustet, 1862), pp.~350--351, \latin{Dominica XVIII. post Pentecosten}.
\item \work{The Roman Missal, Translated into the English Language for the Use of the Laity} (Philadelphia: Eugene Cummiskey, 1861), pp.~444--445.
\item The Holy Bible, Douay--Rheims version, revised by Bishop Richard Challoner: Ex 24; Ps 95; 101; 121; Ecclus 36; 50:29; Mt 8:28--9:9; 1 Cor 1:1--13; 4:7; 16:8. \work{Biblia Sacra Vulgatae editionis} (Clementine Vulgate): Ex 24; 29:41; Lev 8:11; Num 7:1; Ps 95; 101; 121; 140:2; Ecclus 36:18; Mt 9:1--8. \work{The Septuagint} (CATSS text) at Sir 36:15, for the comparison of wording.
\item H.~A. Wilson, ed., \work{The Gelasian Sacramentary} (Oxford, 1894), p.~232 (Liber III, sect.~XIV) and Liber I, sect.~XCIII, and \work{The Gregorian Sacramentary under Charles the Great}, HBS 49 (1915), pp.~174--175.
\item R.-J. Hesbert, \work{Antiphonale Missarum Sextuplex}, as reported by the gregorien.info chant database: \latin{Dominica XVIII post Pentecosten}, and the Gradual \latin{Timebunt gentes}.
\item Germain Morin, ``Le plus ancien \latin{Comes} ou lectionnaire de l'Église romaine'', \work{Revue Bénédictine} 27 (1910), pp.~62--63, and ``Liturgie et basiliques de Rome au milieu du VII\textsuperscript{e} siècle'', \work{Revue Bénédictine} 28 (1911), pp.~315--316 (the Würzburg lists).
\end{itemize}

\subsubsection*{Fathers, Doctors and saints}
\begin{itemize}\setlength{\itemsep}{0.15em}\setlength{\parskip}{0pt}
\item St Augustine, \work{Enarrationes in Psalmos}, Nicene and Post-Nicene Fathers, 1st ser., vol.~8: on Ps~95, 9; Ps~101, s.~1, 16--18; Ps~121, 2, 12. Latin in Migne, PL 37, cols.~1618--1619, 1628. \work{De consensu evangelistarum} II.25.57--58, 1st ser., vol.~6.
\item St Hilary of Poitiers, \work{Tractatus super Psalmos}, ed. A. Zingerle, CSEL 22 (1891), in Ps.~121, 1--5, 14; \work{Commentarius in Matthaeum} VIII.4--8 (Migne, PL 9, cols.~959--962).
\item St Ambrose, \work{Expositio Evangelii secundum Lucam} V.11--15 (Migne, PL 15, cols.~1357--1359).
\item St Bede, \work{In Lucae evangelium expositio} II, on Lk 5:24 (Migne, PL 92, cols.~388--389).
\item St Jerome, \work{Commentarii in Matthaeum} I, on 9:1--8 (Migne, PL 26, cols.~54--56).
\item St John Chrysostom, \work{Homilies on Matthew} 29, Nicene and Post-Nicene Fathers, 1st ser., vol.~10; \work{Homilies on First Corinthians} 2 and \work{Homilies on Hebrews} 16, 1st ser., vols.~12 and 14 (New Advent); \work{Expositio in Psalmum CXXI} (Migne, PG 55, cols.~347--351).
\item Cassiodorus, \work{Expositio psalmorum}, in Ps.~95, 101 and 121 (Migne, PL 70).
\item Theodoret of Cyrus, \work{Interpretatio in Psalmos}, in Ps.~95, 101 and 121 (Migne, PG 80), and \work{Interpretatio in I ad Corinthios}, on 1:4--8 (PG 82, col.~229).
\item \work{Expositiones in Psalmos} printed under the name of St Athanasius, on Ps~95 (Migne, PG 27, cols.~415--416).
\item Ambrosiaster, \work{Commentaria in Epistolam ad Corinthios Primam}, on 1:4--8 (Migne, PL 17).
\item St Thomas Aquinas, \work{Super Evangelium S.~Matthaei lectura} c.~9 (Venice, 1745), pp.~120--122; \work{Super I ad Corinthios} c.~1, lect.~1 (Corpus Thomisticum).
\item St Robert Bellarmine, \work{A Commentary on the Book of Psalms}, trans. J. O'Sullivan (1866, abridged): on Pss~95:8; 121.
\item Cornelius a Lapide, \work{Commentaria in Pentateuchum} (Antwerp, 1700), on Ex 24:4--8; \work{Commentaria in omnes D.~Pauli epistolas} (Antwerp, 1614), on 1 Cor 1:8.
\end{itemize}

\subsubsection*{Liturgical commentators and chronology sources}
\begin{itemize}\setlength{\itemsep}{0.15em}\setlength{\parskip}{0pt}
\item Bl.\ Ildefonso Schuster, \work{The Sacramentary (Liber Sacramentorum)}, vol.~3 (London, 1927), pp.~167--170.
\item The continuation of Prosper Guéranger, \work{The Liturgical Year} (Dom Lucien Fromage), vol.~11, trans. Laurence Shepherd (London, 1909), pp.~393--409.
\item Berno of Reichenau, \work{Libellus de quibusdam rebus ad Missae officium pertinentibus} V (Migne, PL 142, col.~1070).
\item \work{The Catholic Encyclopedia} (New York, 1907--1912): ``Biblical Chronology'' and ``General Chronology'' (1908); ``Epistles to the Corinthians'' (1908); ``Ecclesiasticus'' (1909); ``Gospel of St.\ Matthew'' (1911); ``St.\ Paul'' (1911); ``The New Testament'' and ``Temple of Jerusalem'' (1912).
\item United States Conference of Catholic Bishops, \work{New American Bible, Revised Edition}, introduction to the Psalms, on their dating; summarized, not reproduced.
\item Fraternité Saint-Pierre (France), \work{Ordo du mois}, and Institut du Christ Roi Souverain Prêtre (France), \work{Ordo}, entries for 27~September 2026.
\end{itemize}
\par}
